\clearpage
\phantomsection\addcontentsline{toc}{section}{Soluzioni degli esercizi}
\markboth{Soluzioni degli esercizi}{Soluzioni degli esercizi}
{\headfont\bfseries\Huge\color{navy}Soluzioni degli esercizi\par}
\vspace{4pt}{\color{teal}\rule{3cm}{2.5pt}}\vspace{8pt}

Le soluzioni sono state scritte per questo manuale: per gli esercizi di modellazione sono possibili anche altre risposte corrette, purché giustificate dai requisiti.

\subsection*{Capitolo 1}
\soluzione{1.1}{(a) Logico: dipende dal DBMS. (b) Concettuale: è pensato per essere letto da tutti. (c) Logico: le chiavi esterne sono un concetto del modello relazionale.}
\soluzione{1.2}{Serve se nella stessa ricevuta possono esserci righe identiche (stessa quantità, descrizione e importo) oppure se interessa l'ordine delle righe: senza \emph{Riga} due righe uguali della stessa ricevuta sarebbero indistinguibili.}

\subsection*{Capitolo 2}
\soluzione{2.1}{\emph{Letto} e \emph{Reparto} sono entità (oggetti con proprietà ed esistenza autonoma). \emph{Ricovero} è un legame fra paziente e letto (o reparto); se lo stesso paziente può essere ricoverato più volte nello stesso reparto, va rappresentato come entità (capitolo 3).}
\soluzione{2.2}{Il nome descrive com'è fatta ciascuna occorrenza (uno studente), non l'insieme: per convenzione si usa il singolare.}

\subsection*{Capitolo 3}
\soluzione{3.1}{Una relationship \emph{Visita} fra Paziente e Medico non basta, perché la coppia (paziente, medico) si ripete. Si usa un'entità \emph{Visita} con attributo \emph{Data}, collegata a Paziente e a Medico da due relationship con $(1,1)$ dal lato di Visita e $(0,N)$ dal lato di Paziente e di Medico. Un possibile identificatore: paziente, medico e data (identificatore esterno, come per l'esame ripetuto).}
\soluzione{3.2}{È una relationship ricorsiva su \emph{Persona}. Il legame è simmetrico, quindi i ruoli non aggiungono significato; si possono comunque indicare due ruoli neutri (Coniuge 1, Coniuge 2) per distinguere i due «posti». Con cardinalità $(0,1)$ su entrambi i ruoli ogni persona ha al più un coniuge. Se interessano anche i matrimoni passati, conviene un'entità \emph{Matrimonio}.}

\subsection*{Capitolo 4}
\soluzione{4.1}{La data di inizio della direzione è una proprietà del legame fra docente e dipartimento: attributo della relationship \emph{Direzione}. Il budget è una proprietà del progetto: attributo dell'entità \emph{Progetto} (come già nello schema).}
\soluzione{4.2}{Due attributi: \emph{Email} con cardinalità $(1,1)$ (sottintesa) ed \emph{EmailSecondaria} con cardinalità $(0,1)$. Una sola cardinalità $(1,2)$ non si scrive con i simboli usati nel corso (massimo 1 o N).}

\subsection*{Capitolo 5}
\soluzione{5.1}{Cliente $(0,N)$, Volo $(0,180)$. Con i soli simboli del corso il massimo 180 si scrive N, e il limite di posti diventa un vincolo da documentare (capitolo 8).}
\soluzione{5.2}{Studente $(0,1)$, Docente $(0,N)$: relationship \hl{uno a molti}.}

\subsection*{Capitolo 6}
\soluzione{6.1}{Tre entità in catena. \emph{Palazzo} identificato da un attributo interno (indirizzo o codice). \emph{Scala} identificata da \emph{Lettera} e dal Palazzo, attraverso una relationship con $(1,1)$ dal lato di Scala e $(1,N)$ dal lato di Palazzo. \emph{Appartamento} identificato da \emph{Interno} e dalla Scala, con $(1,1)$ dal lato di Appartamento. Un appartamento è quindi identificato da indirizzo del palazzo, lettera della scala e numero di interno, come un conto da ABI, CAB e numero.}
\soluzione{6.2}{Perché un identificatore esterno richiede che l'entità esterna sia \emph{una e una sola}: con massimo N un reparto apparterrebbe a più divisioni e non si saprebbe quale usare per identificarlo; con minimo 0 potrebbe non appartenere a nessuna.}
\soluzione{6.3}{No. Il codice basta da solo: un identificatore interno semplice. Aggiungere il corso di laurea renderebbe l'identificatore non minimale.}

\subsection*{Capitolo 7}
\soluzione{7.1}{\emph{Utente} generalizza \emph{Socio} e \emph{Ospite giornaliero}: totale (ogni utente è l'uno o l'altro) ed esclusiva. \emph{Istruttore} è un sottoinsieme di \emph{Socio} (generalizzazione con una sola figlia, parziale).}
\soluzione{7.2}{La generalizzazione deve essere totale (lo è) e sia \emph{Auto} sia \emph{Moto} devono avere un proprio identificatore: ogni veicolo è un'auto o una moto e lì viene identificato.}

\subsection*{Capitolo 8}
\soluzione{8.1}{Sede \textbullet\ Sede dell'azienda \textbullet\ Indirizzo, Città \textbullet\ Città.}
\soluzione{8.2}{Per esempio: «un utente non può avere in prestito contemporaneamente due copie dello stesso volume». Riguarda insieme le relationship Prestito e quella fra Copia e Volume, quindi non si esprime con i costrutti del modello.}

\subsection*{Capitolo 9}
\soluzione{9.1}{\emph{Prestito} diventa un'entità con attributi \emph{DataInizio}, \emph{DataPrevista} e \emph{DataRestituzione} (quest'ultima con cardinalità $(0,1)$, vuota per i prestiti in corso). È collegata a Copia e a Utente con $(1,1)$ dal lato di Prestito e $(0,N)$ dal lato di Copia e di Utente. Identificatore: Copia e DataInizio (una copia non inizia due prestiti lo stesso giorno).}
\soluzione{9.2}{\emph{Edizione} ottiene l'identificatore interno \emph{ISBN}. Le relationship verso Collana e Opera restano, con $(1,1)$, ma non fanno più parte dell'identificatore. Volume resta identificato da Numero ed Edizione.}
